% ENEM 2009-2022 — 16 questões MCQ — Regra de Três Simples e Composta
% Fonte: questões selecionadas de provas ENEM (regular e PPL)
% Nota: Q4 do HTML era duplicata de Q3 — omitida

---
tipo: Múltipla Escolha
dificuldade: Média
nivel: médio
disciplina: Matemática
assunto: Regra de Três
gabarito: B
tags: [ENEM, regra de três simples, energia elétrica, bandeiras tarifárias, porcentagem]
fonte: concurso
concurso: ENEM
banca: INEP
ano: 2022
cargo: PPL
numero: 1
---
\question
A tarifa da energia elétrica no Brasil tem sofrido variações em função do seu custo de produção, seguindo um sistema de bandeiras tarifárias. Esse sistema indica se haverá ou não acréscimo no valor do quilowatt-hora (kWh). Suponha que o repasse ao consumidor final seja da seguinte maneira:

\begin{itemize}
\item bandeira verde: a tarifa não sofre acréscimo;
\item bandeira amarela: a tarifa sofre acréscimo de R\$ 0,015 para cada kWh consumido;
\item bandeira vermelha --- patamar 1: a tarifa sofre acréscimo de R\$ 0,04 para cada kWh consumido;
\item bandeira vermelha --- patamar 2: a tarifa sofre acréscimo de R\$ 0,06 para cada kWh consumido.
\end{itemize}

A conta de energia elétrica em uma residência é constituída apenas por um valor correspondente à quantidade de energia elétrica consumida no período medido, multiplicada pela tarifa correspondente. O valor da tarifa em um período com uso da bandeira verde é R\$ 0,42 por kWh consumido. Uma forte estiagem justificou a alteração da bandeira verde para a bandeira vermelha --- patamar 2.

Um usuário, cujo consumo é tarifado na bandeira verde, observa o seu consumo médio mensal. Para não afetar o seu orçamento familiar, ele pretende alterar a sua prática de uso de energia, reduzindo o seu consumo, de maneira que a sua próxima fatura tenha, no máximo, o mesmo valor da conta de energia do período em que era aplicada a bandeira verde.

Qual percentual mínimo de redução de consumo esse usuário deverá praticar de forma a atingir seu objetivo?
\begin{choices}
\choice 6,0\%
\CorrectChoice 12,5\%
\choice 14,3\%
\choice 16,6\%
\choice 87,5\%
\end{choices}

---
tipo: Múltipla Escolha
dificuldade: Média
nivel: médio
disciplina: Matemática
assunto: Regra de Três
gabarito: C
tags: [ENEM, regra de três simples, porcentagem, emprego, Caged]
fonte: concurso
concurso: ENEM
banca: INEP
ano: 2022
numero: 2
---
\question
Em janeiro de 2013, foram declaradas 1.794.272 admissões e 1.765.372 desligamentos no Brasil, ou seja, foram criadas 28.900 vagas de emprego, segundo dados do Cadastro Geral de Empregados e Desempregados (Caged), divulgados pelo Ministério do Trabalho e Emprego (MTE). Segundo o Caged, o número de vagas criadas em janeiro de 2013 sofreu uma queda de 75\%, quando comparado com o mesmo período de 2012.

De acordo com as informações dadas, o número de vagas criadas em janeiro de 2012 foi
\begin{choices}
\choice 16.514
\choice 86.700
\CorrectChoice 115.600
\choice 441.343
\choice 448.568
\end{choices}

---
tipo: Múltipla Escolha
dificuldade: Média
nivel: médio
disciplina: Matemática
assunto: Regra de Três
gabarito: C
tags: [ENEM, regra de três simples, hidratação, atividade física]
fonte: concurso
concurso: ENEM
banca: INEP
ano: 2022
numero: 3
---
\question
Um médico faz o acompanhamento clínico de um grupo de pessoas que realizam atividades físicas diariamente. Ele observou que a perda média de massa dessas pessoas para cada hora de atividade física era de 1,5 kg. Sabendo que a massa de 1 L de água é de 1 kg, ele recomendou que ingerissem, ao longo das 3 horas seguintes ao final da atividade, uma quantidade total de água correspondente a 40\% a mais do que a massa perdida na atividade física, para evitar desidratação.

Seguindo a recomendação médica, uma dessas pessoas ingeriu, certo dia, um total de 1,7 L de água após terminar seus exercícios físicos.

Para que a recomendação médica tenha efetivamente sido respeitada, a atividade física dessa pessoa, nesse dia, durou
\begin{choices}
\choice 30 minutos ou menos.
\choice mais de 35 e menos de 45 minutos.
\CorrectChoice mais de 45 e menos de 55 minutos.
\choice mais de 60 e menos de 70 minutos.
\choice 70 minutos ou mais.
\end{choices}

---
tipo: Múltipla Escolha
dificuldade: Média
nivel: médio
disciplina: Matemática
assunto: Regra de Três
gabarito: D
tags: [ENEM, regra de três simples, desempenho, combustível, km por litro]
fonte: concurso
concurso: ENEM
banca: INEP
ano: 2021
numero: 4
---
\question
Um automóvel apresenta um desempenho médio de 16 km/L. Um engenheiro desenvolveu um novo motor a combustão que economiza, em relação ao consumo do motor anterior, 0,1 L de combustível a cada 20 km percorridos.

O valor do desempenho médio do automóvel com o novo motor, em quilômetro por litro, expresso com uma casa decimal, é
\begin{choices}
\choice 15,9.
\choice 16,1.
\choice 16,4.
\CorrectChoice 17,4.
\choice 18,0.
\end{choices}

---
tipo: Múltipla Escolha
dificuldade: Média
nivel: médio
disciplina: Matemática
assunto: Regra de Três
gabarito: A
tags: [ENEM, regra de três composta, inversamente proporcional, Fórmula 1, pneus]
fonte: concurso
concurso: ENEM
banca: INEP
ano: 2021
cargo: PPL
numero: 5
---
\question
Em uma corrida automobilística, os carros podem fazer paradas nos boxes para efetuar trocas de pneus. Nessas trocas, o trabalho é feito por um grupo de três pessoas em cada pneu. Considere que os grupos iniciam o trabalho no mesmo instante, trabalham à mesma velocidade e cada grupo trabalha em um único pneu. Com os quatro grupos completos, são necessários 4 segundos para que a troca seja efetuada. O tempo gasto por um grupo para trocar um pneu é inversamente proporcional ao número de pessoas trabalhando nele. Em uma dessas paradas, um dos trabalhadores passou mal, não pôde participar da troca e nem foi substituído, de forma que um dos quatro grupos de troca ficou reduzido.

Nessa parada específica, com um dos grupos reduzido, qual foi o tempo gasto, em segundo, para trocar os quatro pneus?
\begin{choices}
\CorrectChoice 6,0
\choice 5,7
\choice 5,0
\choice 4,5
\choice 4,4
\end{choices}

---
tipo: Múltipla Escolha
dificuldade: Média
nivel: médio
disciplina: Matemática
assunto: Regra de Três
gabarito: A
tags: [ENEM, regra de três, planejamento, vendas, desconto]
fonte: concurso
concurso: ENEM
banca: INEP
ano: 2022
cargo: PPL
numero: 6
---
\question
A meta de uma concessionária de automóveis é vender, pelo menos, 104 carros por mês. Sabe-se que, em média, em dias em que não são oferecidos descontos, são vendidos 3 carros por dia; em dias em que há o desconto mínimo, são vendidos 4 carros por dia; e, em dias em que há o desconto máximo, são vendidos 5 carros por dia.

No mês atual, até o fim do expediente do sexto dia em que a concessionária abriu, não foram oferecidos descontos, tendo sido vendidos 18 carros, conforme indicava a média. Ela ainda abrirá por mais 20 dias neste mês.

A menor quantidade de dias em que será necessário oferecer o desconto máximo, de modo que ainda seja possível a concessionária alcançar sua meta de vendas para o mês, é
\begin{choices}
\CorrectChoice 6.
\choice 10.
\choice 11.
\choice 13.
\choice 18.
\end{choices}

---
tipo: Múltipla Escolha
dificuldade: Média
nivel: médio
disciplina: Matemática
assunto: Regra de Três
gabarito: C
tags: [ENEM, regra de três composta, capacidade de carga, materiais de construção]
fonte: concurso
concurso: ENEM
banca: INEP
ano: 2018
cargo: PPL
numero: 7
---
\question
Usando a capacidade máxima de carga do caminhão de uma loja de materiais de construção, é possível levar 60 sacos de cimento, ou 90 sacos de cal, ou 120 latas de areia. No pedido de um cliente, foi solicitada a entrega de 15 sacos de cimento, 30 sacos de cal e a maior quantidade de latas de areia que fosse possível transportar, atingindo a capacidade máxima de carga do caminhão.

Nessas condições, qual a quantidade máxima de latas de areia que poderão ser enviadas ao cliente?
\begin{choices}
\choice 30
\choice 40
\CorrectChoice 50
\choice 80
\choice 90
\end{choices}

---
tipo: Múltipla Escolha
dificuldade: Difícil
nivel: médio
disciplina: Matemática
assunto: Regra de Três
gabarito: D
tags: [ENEM, regra de três, taxa de variação, piscina, escoamento]
fonte: concurso
concurso: ENEM
banca: INEP
ano: 2017
numero: 8
---
\question
Às 17h15min começa uma forte chuva, que cai com intensidade constante. Uma piscina em forma de um paralelepípedo retângulo, que se encontrava inicialmente vazia, começa a acumular a água da chuva e, às 18 horas, o nível da água em seu interior alcança 20 cm de altura. Nesse instante, é aberto o registro que libera o escoamento da água por um ralo localizado no fundo dessa piscina, cuja vazão é constante. Às 18h40min a chuva cessa e, nesse exato instante, o nível da água na piscina baixou para 15 cm.

O instante em que a água dessa piscina terminar de escoar completamente está compreendido entre
\begin{choices}
\choice 19h30min e 20h10min
\choice 19h20min e 19h30min
\choice 19h10min e 19h20min
\CorrectChoice 19h e 19h10min
\choice 18h40min e 19h
\end{choices}

---
tipo: Múltipla Escolha
dificuldade: Média
nivel: médio
disciplina: Matemática
assunto: Regra de Três
gabarito: B
tags: [ENEM, regra de três composta, banco de sangue, plasma, saúde]
fonte: concurso
concurso: ENEM
banca: INEP
ano: 2016
cargo: PPL
numero: 9
---
\question
Um banco de sangue recebe 450 mL de sangue de cada doador. Após separar o plasma sanguíneo das hemácias, o primeiro é armazenado em bolsas de 250 mL de capacidade. O banco de sangue aluga refrigeradores de uma empresa para estocagem das bolsas de plasma, segundo a sua necessidade. Cada refrigerador tem uma capacidade de estocagem de 50 bolsas. Ao longo de uma semana, 100 pessoas doaram sangue àquele banco.

Admita que, de cada 60 mL de sangue, extraem-se 40 mL de plasma.

O número mínimo de congeladores que o banco precisou alugar, para estocar todas as bolsas de plasma dessa semana, foi
\begin{choices}
\choice 2.
\CorrectChoice 3.
\choice 4.
\choice 6.
\choice 8.
\end{choices}

---
tipo: Múltipla Escolha
dificuldade: Fácil
nivel: médio
disciplina: Matemática
assunto: Regra de Três
gabarito: C
tags: [ENEM, regra de três simples, dosagem, veterinária, antibiótico]
fonte: concurso
concurso: ENEM
banca: INEP
ano: 2015
cargo: PPL
numero: 10
---
\question
Um granjeiro detectou uma infecção bacteriológica em sua criação de 100 coelhos. A massa de cada coelho era de, aproximadamente, 4 kg. Um veterinário prescreveu a aplicação de um antibiótico, vendido em frascos contendo 16 mL, 25 mL, 100 mL, 400 mL ou 1.600 mL. A bula do antibiótico recomenda que, em aves e coelhos, seja administrada uma dose única de 0,25 mL para cada quilograma de massa do animal.

Para que todos os coelhos recebessem a dosagem do antibiótico recomendada pela bula, de tal maneira que não sobrasse produto na embalagem, o criador deveria comprar um único frasco com a quantidade, em mililitros, igual a
\begin{choices}
\choice 16.
\choice 25.
\CorrectChoice 100.
\choice 400.
\choice 1.600.
\end{choices}

---
tipo: Múltipla Escolha
dificuldade: Média
nivel: médio
disciplina: Matemática
assunto: Regra de Três
gabarito: C
tags: [ENEM, regra de três, câmbio, moedas, dólar, euro, libra]
fonte: concurso
concurso: ENEM
banca: INEP
ano: 2013
cargo: PPL
numero: 11
---
\question
A cotação de uma moeda em relação a uma segunda moeda é o valor que custa para comprar uma unidade da primeira moeda, utilizando a segunda moeda. Por exemplo, se a cotação do dólar é 1,6 real, isso significa que para comprar 1 dólar é necessário 1,6 real.

Suponha que a cotação do dólar, em reais, seja de 1,6 real, a do euro, em reais, seja de 2,4 reais e a cotação da libra, em euros, seja de 1,1 euro.

Qual é a cotação da libra, em dólares?
\begin{choices}
\choice 4,224 dólares
\choice 2,64 dólares
\CorrectChoice 1,65 dólar
\choice 1,50 dólar
\choice 1,36 dólar
\end{choices}

---
tipo: Múltipla Escolha
dificuldade: Média
nivel: médio
disciplina: Matemática
assunto: Regra de Três
gabarito: C
tags: [ENEM, regra de três composta, reservatório, ralos, escoamento]
fonte: concurso
concurso: ENEM
banca: INEP
ano: 2013
numero: 12
---
\question
Uma indústria tem um reservatório de água com capacidade para 900 m$^3$. Quando há necessidade de limpeza do reservatório, toda a água precisa ser escoada. O escoamento da água é feito por seis ralos, e dura 6 horas quando o reservatório está cheio. Esta indústria construirá um novo reservatório, com capacidade de 500 m$^3$, cujo escoamento da água deverá ser realizado em 4 horas, quando o reservatório estiver cheio. Os ralos utilizados no novo reservatório deverão ser idênticos aos do já existente.

A quantidade de ralos do novo reservatório deverá ser igual a
\begin{choices}
\choice 2.
\choice 4.
\CorrectChoice 5.
\choice 8.
\choice 9.
\end{choices}

---
tipo: Múltipla Escolha
dificuldade: Fácil
nivel: médio
disciplina: Matemática
assunto: Regra de Três
gabarito: D
tags: [ENEM, regra de três simples, shopping, tíquetes]
fonte: concurso
concurso: ENEM
banca: INEP
ano: 2012
numero: 13
---
\question
Nos \textit{shopping centers} costumam existir parques com vários brinquedos e jogos. Os usuários colocam créditos em um cartão, que são descontados por cada período de tempo de uso dos jogos. Dependendo da pontuação da criança no jogo, ela recebe um certo número de tíquetes para trocar por produtos nas lojas dos parques. Suponha que o período de uso de um brinquedo em certo \textit{shopping} custa R\$ 3,00 e que uma bicicleta custa 9.200 tíquetes.

Para uma criança que recebe 20 tíquetes por período de tempo que joga, o valor, em reais, gasto com créditos para obter a quantidade de tíquetes para trocar pela bicicleta é
\begin{choices}
\choice 153.
\choice 460.
\choice 1.218.
\CorrectChoice 1.380.
\choice 3.066.
\end{choices}

---
tipo: Múltipla Escolha
dificuldade: Fácil
nivel: médio
disciplina: Matemática
assunto: Regra de Três
gabarito: B
tags: [ENEM, regra de três simples, bacia sanitária, economia de água, ecologia]
fonte: concurso
concurso: ENEM
banca: INEP
ano: 2012
numero: 14
---
\question
Há, em virtude da demanda crescente de economia de água, equipamentos e utensílios como, por exemplo, as bacias sanitárias ecológicas, que utilizam 6 litros de água por descarga em vez dos 15 litros utilizados por bacias sanitárias não ecológicas, conforme dados da Associação Brasileira de Normas Técnicas (ABNT).

Qual será a economia diária de água obtida por meio da substituição de uma bacia sanitária não ecológica, que gasta cerca de 60 litros por dia com a descarga, por uma bacia sanitária ecológica?
\begin{choices}
\choice 24 litros
\CorrectChoice 36 litros
\choice 40 litros
\choice 42 litros
\choice 50 litros
\end{choices}

---
tipo: Múltipla Escolha
dificuldade: Difícil
nivel: médio
disciplina: Matemática
assunto: Regra de Três
gabarito: D
tags: [ENEM, regra de três composta, trabalhadores, máquinas, colheita]
fonte: concurso
concurso: ENEM
banca: INEP
ano: 2009
numero: 15
---
\question
Uma cooperativa de colheita propôs a um fazendeiro um contrato de trabalho nos seguintes termos: a cooperativa forneceria 12 trabalhadores e 4 máquinas, em um regime de trabalho de 6 horas diárias, capazes de colher 20 hectares de milho por dia, ao custo de R\$ 10,00 por trabalhador por dia de trabalho, e R\$ 1.000,00 pelo aluguel diário de cada máquina. O fazendeiro argumentou que fecharia contrato se a cooperativa colhesse 180 hectares de milho em 6 dias, com gasto inferior a R\$ 25.000,00.

Para atender às exigências do fazendeiro e supondo que o ritmo dos trabalhadores e das máquinas seja constante, a cooperativa deveria
\begin{choices}
\choice manter sua proposta.
\choice oferecer 4 máquinas a mais.
\choice oferecer 6 trabalhadores a mais.
\CorrectChoice aumentar a jornada de trabalho para 9 horas diárias.
\choice reduzir em R\$ 400,00 o valor do aluguel diário de uma máquina.
\end{choices}

---
tipo: Múltipla Escolha
dificuldade: Média
nivel: médio
disciplina: Matemática
assunto: Regra de Três
gabarito: B
tags: [ENEM, regra de três composta, Fórmula 1, combustível, densidade]
fonte: concurso
concurso: ENEM
banca: INEP
ano: 2009
numero: 16
---
\question
Segundo as regras da Fórmula 1, o peso mínimo do carro, de tanque vazio, com o piloto, é de 605 kg, e a gasolina deve ter densidade entre 725 e 780 gramas por litro. Entre os circuitos nos quais ocorrem competições dessa categoria, o mais longo é Spa-Francorchamps, na Bélgica, cujo traçado tem 7 km de extensão. O consumo médio de um carro da Fórmula 1 é de 75 litros para cada 100 km.

Suponha que um piloto de uma equipe específica, que utiliza um tipo de gasolina com densidade de 750 g/L, esteja no circuito de Spa-Francorchamps, parado no box para reabastecimento. Caso ele pretenda dar mais 16 voltas, ao ser liberado para retornar à pista, seu carro deverá pesar, no mínimo,
\begin{choices}
\choice 617 kg.
\CorrectChoice 668 kg.
\choice 680 kg.
\choice 689 kg.
\choice 717 kg.
\end{choices}
